\chapter{Missing data, structural zeros, and what is not implemented}\label{ch:missing}

The first part of this chapter explains how \nwsaom{} treats data that are unknown (missing ties and behavior values) or fixed by design (structural zeros and ones). The second part is an honest list of what RSiena offers and \nwsaom{} does not.

\section{Missing data: \opt{missnet()} and \opt{missbeh()}}\label{sec:missing}
Missing data here means dyads or actors whose true value at a wave is \emph{unknown}. This is RSiena's regular missing-data mechanism and it is different from composition change (actors who are not part of the network at a wave, \opt{present()} or \opt{joinat()}/\opt{leaveat()}): a missing tie is uncertainty about an otherwise active actor. The two can be used together.

\paragraph{Specifying missingness.}
\begin{itemize}
\item \opt{missnet(matlist)} takes one 0/1 $n\times n$ \emph{matrix} name per wave, in the order of \opt{wave1()}/\opt{wave2()} or \opt{waves()}; a 1 marks a dyad missing at that wave. It is a raw Stata matrix, not an \texttt{nwset} network. Build it with \texttt{matrix input} or \texttt{mkmat} from your own indicator. With two waves, \texttt{missnet(m1 m2)}.
\item \opt{missbeh(varlist)} takes one 0/1 \emph{variable} per wave (the same one-variable-per-wave convention as \opt{present()}), marking the actors whose behavior value is missing at that wave. It requires \opt{behavior()}.
\item Both are optional and independent: use either or both.
\end{itemize}
\begin{lstlisting}
nwwebuse glasgow, nwclear
* the first five actors did not report their ties at wave 2
matrix m1 = J(50,50,0)
matrix m2 = J(50,50,0)
forvalues i = 1/5 {
    forvalues j = 1/50 {
        if `i' != `j' matrix m2[`i',`j'] = 1
    }
}
nwsaom, wave1(glasgow1) wave2(glasgow2) outdegree reciprocity transtrip ///
    missnet(m1 m2) seed(1)
\end{lstlisting}
(This example was run with the development build; it converged, with an overall convergence ratio of 0.119.) Because the missing indicators are Stata matrices, they are limited by Stata's matrix size: 800$\times$800 in Stata/BE, 11,000$\times$11,000 in Stata/SE.

\paragraph{How missing values are handled.} \nwsaom{} follows RSiena's two steps.
\begin{enumerate}
\item \textbf{Imputation.} Each missing dyad is filled in by carrying forward the last observed value (0 if never observed up to that wave). A missing behavior value is filled in from the previous observation, else the next observation, else the mode of that wave (cross-sectional). The imputed values take part in the simulation completely normally; unlike \opt{present()}, missingness does not restrict who can act.
\item \textbf{Masking.} A dyad or actor missing at \emph{either} end wave of a period is excluded from the observed target statistic \emph{and} from every simulated statistic of that period, so the moment condition is not biased by the exclusion. The same applies to the distance that conditional estimation simulates to, and to the distance statistic of the rate under unconditional estimation. The existing effects needed no special-casing: their statistic functions are reused unchanged.
\end{enumerate}
Missing data are handled in the native engine, when the model is otherwise native-eligible; the masked statistic matches the Mata engine to machine precision.

\paragraph{A disclosed approximation.} The statistic of \opt{avalt} and \opt{avsim} for an actor depends on the \emph{current} values of the actor's alters, which differ across simulated replicates. For these effects the masking is only an approximation (it is exact for \opt{linear} and \opt{quadratic}). It is empirically better than not masking, but may be less precise than an idealised per-effect masking under heavy missingness combined with \opt{avalt} or \opt{avsim}.

\paragraph{Missing covariate values.} Covariates of \opt{nodematch()}, \opt{nodecov()}, \opt{nodeicov()}, \opt{nodeocov()}, \opt{simcov()} and their aliases follow RSiena. A missing value is imputed with the covariate mean when actors decide their ministeps, but the \emph{statistic} (target and every simulated statistic) skips a tie whose relevant actor is missing: the alter for \opt{altx()}, the ego for \opt{egox()}, either for \opt{samex()} and \opt{simx()}. With the default centring a missing \opt{altx()}, \opt{egox()} or \opt{nodecov()} value becomes 0 and skipping changes nothing; under \opt{nocenter} and for \opt{samex()}/\opt{simx()} it does, and the range and similarity mean of \opt{simx()} use the observed values only. On RSiena's \texttt{s501}/\texttt{s502} with 8 of 50 covariate values missing, density, reciprocity, altX, egoX and simX agree with RSiena 1.6.6 within sampling noise (simX $0.79$ against $0.78$). An \opt{interact()} naming a \opt{samex()} or \opt{simx()} effect whose variable has missing values is rejected.

\section{Structural zeros and ones}\label{sec:structural}
A structural dyad is one whose tie value is known and cannot change, for example a legally mandated reporting tie, a physically impossible tie, or a dyad that you wish to hold fixed for a counterfactual. It is a simpler idea than missing data: the value is not unknown. \nwsaom{} follows RSiena: the values are part of the wave networks. In a wave network a tie with the value 10 is a structural zero and a tie with the value 11 a structural one (a 1 is an ordinary tie, a 0 no tie). \nwsaom{} detects the codes; there is no option. With \texttt{nwset, mat()} the codes are the cell values of the matrix, with \texttt{nwfromedge} give the value variable the values 10 and 11. Codes on the diagonal are ignored; every other tie value than 1, 10 and 11 stops the command.

The rules are RSiena's, checked against RSiena 1.6.14 on the data of the examples:
\begin{itemize}
\item The codes belong to a wave; the structural dyads may differ from wave to wave. A dyad that is structural at the \emph{start} of a period keeps its start value during the whole period (the value of its code: 11 is a tie, 10 no tie) and is not among the alternatives of the actors. A dyad that is structural at the \emph{end} of a period only is free during the period and gets its end value before the statistics are computed.
\item The observed statistics of a period are those of the end network with the dyads that were structural at the start of the period set to their start value. The observed distance between two waves, which is the target of the rate and of the conditional estimation, leaves out the dyads that are structural at either end of the period and the dyads that are missing at either end. The conditional estimation stops, as RSiena's, when the number of changes of the dyads that are not missing reaches that distance, so with structural dyads that change from wave to wave the unconditional estimation (option \opt{unconditional}) is usually the better choice.
\item An actor with no free alternative among the active actors has the rate 0; the other actors cannot choose a structural dyad.
\item A dyad that is missing at a wave (\opt{missnet()}) takes the code of the previous wave. Composition change (\opt{present()}, \opt{joinat()}, \opt{leaveat()}), covariates, interactions, rate covariates and the structural rate effects (computed over the actors who can act), network endowment and creation (RSiena advises against them when the structural dyads change from wave to wave) and two-mode networks (the codes are cells of the actors-by-events matrix) work as without codes.
\end{itemize}
The following example makes the ties from pupil 1 to the pupils 11 to 20 structural zeros in both waves:
\begin{lstlisting}
nwtomata glasgow1, mat(G1)
nwtomata glasgow2, mat(G2)
forvalues j = 11/20 {
    mata: G1[1,`j'] = 10; G2[1,`j'] = 10
}
mata: st_matrix("g1m", G1); st_matrix("g2m", G2)
nwset, mat(g1m) directed name(g1)
nwset, mat(g2m) directed name(g2)
nwsaom, wave1(g1) wave2(g2) outdegree reciprocity transtrip seed(1)
\end{lstlisting}

\paragraph{Scope.} One directed one-mode network or one two-mode network, any number of waves, the method of moments. One co-evolving behavior is possible (without rate covariates). Not yet with several behaviors, \opt{groups()}, \opt{symmetric}, \texttt{nwsaom multinet}, \opt{ml} and \texttt{nwsaom bayes}; \texttt{estat gof}, \texttt{mems}, \texttt{simulate}, \texttt{importance}, \texttt{score}, \texttt{timetest} and \texttt{nwsaom describe} stop with a message after or for such data. The code runs in the native plugin only (there is no Mata implementation).

\paragraph{The older option \opt{structural(matname)}.} It takes one 0/1 matrix of dyads that are fixed in both periods and need the same observed value at both waves; two waves, network only. It is kept for compatibility, is replaced by the codes, and cannot be combined with them.

\section{What RSiena has that \nwsaom{} does not (yet)}\label{sec:notimpl}
RSiena's effects table, as returned by \texttt{getEffects()} for twelve data configurations of RSiena 1.6.6, has \nTotal{} rows (a row is one RSiena short name for one variable type and interaction context). The effects that RSiena 1.6.7 to 1.6.14 added or changed (for example \texttt{WXZX} and the two-mode \texttt{toU}) are implemented and compared with RSiena 1.6.14 but are not part of these counts; the counts change with RSiena releases. Every row was compared with \nwsaom{} by reading the code and the documentation, and by running models against \texttt{siena07} where they were implemented. Of these rows:
\begin{itemize}
\item \nImpl{} are implemented, \nPartial{} is partly implemented;
\item \nNotImpl{} are \emph{not implemented};
\item \nNotPlanned{} are deliberately out of scope: continuous (SDE) behavior and the effects that involve a continuous dependent variable.
\end{itemize}
Of the \nPlannable{} rows that are not out of scope, \pctImpl\% are implemented. The percentage is by \emph{rows}, not by how often an effect is used: the effects most often used in applied papers are largely covered, and the missing rows are concentrated in a few groups. The groups follow; the row-level list is at the end of the chapter.

\subsection{Effects not implemented, by reason}
The \nImpl{} implemented rows are \pctImpl\% of the plannable ones. The \nNotImpl{} rows that are not implemented are of five kinds.
\begin{description}[leftmargin=0pt,labelindent=0pt]
\item[The generalised method of moments (31 rows).] RSiena's GMoM estimator uses its own statistics (the effects named like \texttt{egoX\_gmm}, \texttt{avAlt\_gmm}, \texttt{crprod\_gmm}, and related placeholders). \nwsaom{} implements the standard method of moments, and maximum likelihood and Bayesian estimation, but not the generalised method of moments.
\item[The settings model (12 rows).] The \texttt{setting*} effects (\texttt{settingSizeAct}, \texttt{settingOppAct} and relatives) belong to RSiena's settings model, a separate model family with setting steps and a rate for each setting; \nwsaom{} has the standard tie-change model only. There is also nothing to validate against by estimation: in \texttt{siena07} eight of the twelve effects have a statistic that is constant under simulation (the estimate stays 0 and has no standard error), and the four variants with the suffix D have a zero derivative and are fixed at 0.
\item[Effects that RSiena cannot evaluate or does not offer (15 rows).] \texttt{siena07} stops or returns nothing useful for these combinations, so there is nothing to compare with and \nwsaom{} refuses them with an error message (Table below and Section~\ref{sec:cannoteval}).
\item[Placeholders (2 rows).] \texttt{minXAlt} and \texttt{maxXAlt} in a non-directed network with unspecified variables are entries of RSiena's effects table that \texttt{siena07} does not offer.
\item[Two further effects.] \texttt{sameWXClosure} for a two-mode network together with a co-evolving behavior inside \texttt{nwsaom multinet} (a two-mode member and a behavior cannot be combined there), and \texttt{RateX} for a continuous behavior variable (continuous behavior is not planned).
\end{description}
\noindent In addition \nNotPlanned{} rows are out of scope: continuous (stochastic differential equation) dependent variables and all effects that involve them.

\subsection{Effects that RSiena cannot evaluate}\label{sec:cannoteval}
\begin{center}\small
\begin{tabular}{@{}p{4.2cm}p{10.4cm}@{}}
\toprule
Effect & What \texttt{siena07} does\\
\midrule
\texttt{inPop\_ego}, two-mode network & stops with a C++ exception\\
\texttt{sameXOutAct}, \texttt{diffXOutAct}, two-mode network, with a constant covariate, a changing covariate or a behavior & the R process aborts\\
\texttt{avGroupEgoX}, \texttt{totGroupEgoX}, two-mode network, constant covariate & reports the effect as not meaningful (available with a changing covariate)\\
\texttt{homXOutAct2} with a behavior & reports the effect as not implemented for behavioral variables\\
\texttt{inRateInv}, \texttt{inRateLog}, two-mode network & stops with an error (the receivers of ties are not the actors)\\
\texttt{susceptAvCovar}, \texttt{infectCovar} with a behavior as covariate & no such effect exists (the covariate is a constant or changing covariate)\\
\texttt{outOutDist2AvIntn} & listed in the effects table, but there is no implementation\\
\texttt{toU} for a one-mode network that depends on a two-mode one & rejected\\
\texttt{avAlt.2M.tie} with RSiena's default parameter 3 & throws an error; \nwsaom{} offers the parameters 1 and 2\\
\bottomrule
\end{tabular}
\end{center}
Further, \opt{higher} in a non-directed network equals the density statistic and is not identified: RSiena returns a standard error of NA and \nwsaom{} stops.

\subsection{RSiena features beyond effects}\label{sec:est-notimpl}
Besides effects, RSiena provides estimation and analysis tools. This is how they compare.
\begin{description}[leftmargin=0pt,labelindent=0pt]
\item[Method of moments (\texttt{siena07}).] Implemented, conditional and unconditional, with the same phase structure (Chapter~\ref{ch:estimation}).
\item[Maximum likelihood and Bayesian estimation (\texttt{siena07} with \texttt{maxlike = TRUE}, \texttt{multi\_siena}).] Available for a directed one-mode network, with or without one co-evolving behavior: the option \opt{ml} and the subcommand \texttt{nwsaom bayes} (Chapter~\ref{ch:estimation}). Not available for non-directed and two-mode networks, several networks, missing data, composition change, rate covariates and several groups. The hierarchical Bayesian analysis of several groups (\texttt{multi\_siena} with several groups) and Bayesian estimation of a behavior alone are not implemented.
\item[Multi-group analysis (\texttt{siena08}, \texttt{sienaGroupCreate}).] Both ways are available. The meta-analysis of separate fits is \texttt{nwsaom meta} (random-effects pooling by robust IWLS, heterogeneity tests, Fisher combination, maximum-likelihood mean and between-group sd with profile intervals). RSiena's joint multi-group \emph{estimation} (one model fitted to several groups at once) is the option \opt{groups()} (Chapter~\ref{ch:groups}), by the method of moments, conditional and unconditional. Not yet implemented with groups: maximum likelihood and Bayesian estimation, \texttt{estat gof}, \texttt{estat simulate} and \texttt{estat importance} (\texttt{estat score} and \texttt{estat timetest} work), two-mode and non-directed networks, \texttt{nwsaom multinet} and several behaviors, the effects that use a constant pooled over the data (the centered degree effects \opt{inpopc}, \opt{outpopc}, \opt{outactc}, \opt{inactc}, \opt{degplus}, \opt{degplusc}, \opt{avdeg}, \opt{transtriads}, structural equivalence), \opt{behrate()}, \opt{structural()}, \opt{joinat()} and \opt{leaveat()}. The hierarchical Bayesian analysis of several groups is not implemented. With groups, missing behavior values together with \opt{behegox}/\opt{behaltx} differ from RSiena, and leavers were not compared (Section~\ref{sec:differences}).
\item[Goodness of fit (\texttt{sienaGOF}).] Available as \texttt{estat gof} (Chapter~\ref{ch:estimation}), with the differences listed there, and not for fits with several co-evolving behaviors.
\item[Composition change.] Available, including fractional join/leave times (\opt{joinat()}/\opt{leaveat()}) as in \opt{sienaCompositionChange()} (with \opt{groups()} only \opt{present()}). Composition change and missing ties are not available inside \texttt{nwsaom multinet}.
\item[Continuous behavior.] Not implemented (out of scope).
\item[Simulation only, time-heterogeneity tests, score-type tests, effect importance.] Available: \texttt{nwsaom simulate}, \texttt{estat timetest}, \texttt{estat score} and \texttt{estat importance} (Chapter~\ref{ch:tools}). After \texttt{nwsaom multinet} the tests are available, but not \texttt{estat gof}, \texttt{estat importance} and \texttt{estat simulate}.
\item[Fixing and testing individual parameters.] Available: \opt{fix()} holds an effect at its \opt{theta0()} value (RSiena's \texttt{fix}), and \texttt{estat score} tests it, also after \texttt{nwsaom multinet}.
\item[The generalised method of moments.] Not implemented (31 statistics).
\item[The settings model.] Not implemented (12 effects).
\item[Model types for non-directed relations.] Offered by name or RSiena's modelType number (\opt{symtype()}: forcing, confirmation, force, agree, joint; modelType 2, 3, 4, 5, 6).
\end{description}

\section{Row-level list of the effects that are not implemented}\label{sec:rowlist}
The tables below list every RSiena effects-table row that is \emph{not implemented} (one table per reason), in the order of RSiena's effects table; the final entry is the out-of-scope group. A row names an RSiena short name together with the variable type and the interaction context (another network, a constant or changing covariate, a dyadic covariate, a behavior variable). Rows that are implemented are listed in the appendix on effects.

\subsection*{GMoM statistics (the generalised method of moments estimator is not implemented) (31 rows)}
\footnotesize
\begin{longtable}{@{}p{3.6cm}p{4.0cm}p{6.4cm}@{}}
\toprule
RSiena name & Network / context & Note\\
\midrule
\endhead
\opt{newrecip} & gmm (one-mode network) & GMoM statistic: nwsaom estimates by Method of Moments only (no GMoM estimator).\\
\opt{persistrecip} & gmm (one-mode network) & GMoM statistic: nwsaom estimates by Method of Moments only (no GMoM estimator).\\
\opt{realrecip} & gmm (one-mode network) & GMoM statistic: nwsaom estimates by Method of Moments only (no GMoM estimator).\\
\opt{realtrans} & gmm (one-mode network) & GMoM statistic: nwsaom estimates by Method of Moments only (no GMoM estimator).\\
\opt{agreetrans} & gmm (one-mode network) & GMoM statistic: nwsaom estimates by Method of Moments only (no GMoM estimator).\\
\opt{egoX\_gmm} & gmm (one-mode network) + constant covariate & GMoM statistic: nwsaom estimates by Method of Moments only (no GMoM estimator).\\
\opt{simX\_gmm} & gmm (one-mode network) + constant covariate & GMoM statistic: nwsaom estimates by Method of Moments only (no GMoM estimator).\\
\opt{egoX\_gmm} & gmm (one-mode network) + changing covariate & GMoM statistic: nwsaom estimates by Method of Moments only (no GMoM estimator).\\
\opt{simX\_gmm} & gmm (one-mode network) + changing covariate & GMoM statistic: nwsaom estimates by Method of Moments only (no GMoM estimator).\\
\opt{egoX\_gmm} & gmm (one-mode network) + a behavior variable & GMoM statistic: nwsaom estimates by Method of Moments only (no GMoM estimator).\\
\opt{simX\_gmm} & gmm (one-mode network) + a behavior variable & GMoM statistic: nwsaom estimates by Method of Moments only (no GMoM estimator).\\
\opt{avSim\_gmm} & gmm (behavior) + another network & GMoM statistic: nwsaom estimates by Method of Moments only (no GMoM estimator).\\
\opt{avAlt\_gmm} & gmm (behavior) + another network & GMoM statistic: nwsaom estimates by Method of Moments only (no GMoM estimator).\\
\opt{totAlt\_gmm} & gmm (behavior) + another network & GMoM statistic: nwsaom estimates by Method of Moments only (no GMoM estimator).\\
\opt{maxAlt\_gmm} & gmm (behavior) + another network & GMoM statistic: nwsaom estimates by Method of Moments only (no GMoM estimator).\\
\opt{minAlt\_gmm} & gmm (behavior) + another network & GMoM statistic: nwsaom estimates by Method of Moments only (no GMoM estimator).\\
\opt{crprod\_gmm} & gmm (one-mode network) + another network & GMoM statistic: nwsaom estimates by Method of Moments only (no GMoM estimator).\\
\opt{from\_gmm} & gmm (one-mode network) + another network & GMoM statistic: nwsaom estimates by Method of Moments only (no GMoM estimator).\\
\opt{to\_gmm} & gmm (one-mode network) + another network & GMoM statistic: nwsaom estimates by Method of Moments only (no GMoM estimator).\\
\opt{to\_gmm} & gmm (two-mode network) + another network & GMoM statistic: nwsaom estimates by Method of Moments only (no GMoM estimator).\\
\opt{egoX\_gmm} & gmm (one-mode network) + continuous variable & GMoM statistic: nwsaom estimates by Method of Moments only (no GMoM estimator).\\
\opt{simX\_gmm} & gmm (one-mode network) + continuous variable & GMoM statistic: nwsaom estimates by Method of Moments only (no GMoM estimator).\\
\opt{egoX\_gmm} & gmm (non-directed network) + constant covariate & GMoM statistic: nwsaom estimates by Method of Moments only (no GMoM estimator).\\
\opt{simX\_gmm} & gmm (non-directed network) + constant covariate & GMoM statistic: nwsaom estimates by Method of Moments only (no GMoM estimator).\\
\opt{egoX\_gmm} & gmm (non-directed network) + changing covariate & GMoM statistic: nwsaom estimates by Method of Moments only (no GMoM estimator).\\
\opt{simX\_gmm} & gmm (non-directed network) + changing covariate & GMoM statistic: nwsaom estimates by Method of Moments only (no GMoM estimator).\\
\opt{egoX\_gmm} & gmm (non-directed network) + a behavior variable & GMoM statistic: nwsaom estimates by Method of Moments only (no GMoM estimator).\\
\opt{simX\_gmm} & gmm (non-directed network) + a behavior variable & GMoM statistic: nwsaom estimates by Method of Moments only (no GMoM estimator).\\
\opt{from\_gmm} & gmm (non-directed network) + another network & GMoM statistic: nwsaom estimates by Method of Moments only (no GMoM estimator).\\
\opt{to\_gmm} & gmm (non-directed network) + another network & GMoM statistic: nwsaom estimates by Method of Moments only (no GMoM estimator).\\
\opt{crprod\_gmm} & gmm (non-directed network) + another network & GMoM statistic: nwsaom estimates by Method of Moments only (no GMoM estimator).\\
\bottomrule
\end{longtable}
\normalsize

\subsection*{Effects that RSiena cannot evaluate or does not offer (15 rows)}
\footnotesize
\begin{longtable}{@{}p{3.6cm}p{4.0cm}p{6.4cm}@{}}
\toprule
RSiena name & Network / context & Note\\
\midrule
\endhead
\opt{inPop\_ego} & two-mode network & Not implemented: siena07 itself stops with a C++ exception for this effect in a two-mode network, so there is nothing to compare with.\\
\opt{sameXOutAct} & two-mode network + constant covariate & Not implemented: siena07 crashes (R aborts) when it evaluates this effect in a two-mode network, so nwsaom refuses it.\\
\opt{diffXOutAct} & two-mode network + constant covariate & Not implemented: siena07 crashes (R aborts) when it evaluates this effect in a two-mode network, so nwsaom refuses it.\\
\opt{avGroupEgoX} & two-mode network + constant covariate & Not implemented: RSiena reports the effect as not meaningful for a constant covariate.\\
\opt{totGroupEgoX} & two-mode network + constant covariate & Not implemented: RSiena reports the effect as not meaningful for a constant covariate.\\
\opt{sameXOutAct} & two-mode network + changing covariate & Not implemented: siena07 crashes (R aborts) when it evaluates this effect in a two-mode network, so nwsaom refuses it.\\
\opt{diffXOutAct} & two-mode network + changing covariate & Not implemented: siena07 crashes (R aborts) when it evaluates this effect in a two-mode network, so nwsaom refuses it.\\
\opt{homXOutAct2} & one-mode network + a behavior variable & Not implemented: RSiena reports the effect as not implemented for a behavior variable.\\
\opt{sameXOutAct} & two-mode network + a behavior variable & Not implemented: siena07 crashes (R aborts) when it evaluates this effect in a two-mode network, so nwsaom refuses it.\\
\opt{diffXOutAct} & two-mode network + a behavior variable & Not implemented: siena07 crashes (R aborts) when it evaluates this effect in a two-mode network, so nwsaom refuses it.\\
\opt{homXOutAct2} & non-directed network + a behavior variable & Not implemented: RSiena reports the effect as not implemented for a behavior variable.\\
\opt{susceptAvCovar} & rate (behavior) + another network + a behavior variable & Not implemented: RSiena has no such effect (the covariate can only be a constant or a changing covariate).\\
\opt{infectCovar} & rate (behavior) + another network + a behavior variable & Not implemented: RSiena has no such effect (the covariate can only be a constant or a changing covariate).\\
\opt{inRateInv} & rate (two-mode network) & Not implemented: siena07 stops with an error for this effect in a two-mode network (the receivers of ties are not the actors).\\
\opt{inRateLog} & rate (two-mode network) & Not implemented: siena07 stops with an error for this effect in a two-mode network (the receivers of ties are not the actors).\\
\bottomrule
\end{longtable}
\normalsize

\subsection*{RSiena's settings model (a different simulation procedure) (12 rows)}
\footnotesize
\begin{longtable}{@{}p{3.6cm}p{4.0cm}p{6.4cm}@{}}
\toprule
RSiena name & Network / context & Note\\
\midrule
\endhead
\opt{settingSizeAct} & one-mode network & Not implemented: RSiena's settings model is a separate model family (setting steps with rates per setting). In siena07 the effect has no estimate to compare wi…\\
\opt{settingSizeActSqrt} & one-mode network & Not implemented: RSiena's settings model is a separate model family (setting steps with rates per setting). In siena07 the effect has no estimate to compare wi…\\
\opt{settingSizeActLog} & one-mode network & Not implemented: RSiena's settings model is a separate model family (setting steps with rates per setting). In siena07 the effect has no estimate to compare wi…\\
\opt{settingSizeActInv} & one-mode network & Not implemented: RSiena's settings model is a separate model family (setting steps with rates per setting). In siena07 the effect has no estimate to compare wi…\\
\opt{settingOppAct} & one-mode network & Not implemented: RSiena's settings model is a separate model family (setting steps with rates per setting). In siena07 the effect has no estimate to compare wi…\\
\opt{settingOppActSqrt} & one-mode network & Not implemented: RSiena's settings model is a separate model family (setting steps with rates per setting). In siena07 the effect has no estimate to compare wi…\\
\opt{settingOppActLog} & one-mode network & Not implemented: RSiena's settings model is a separate model family (setting steps with rates per setting). In siena07 the effect has no estimate to compare wi…\\
\opt{settingLogCreationAct} & one-mode network & Not implemented: RSiena's settings model is a separate model family (setting steps with rates per setting). In siena07 the effect has no estimate to compare wi…\\
\opt{settingOppActD} & one-mode network & Not implemented: RSiena's settings model is a separate model family (setting steps with rates per setting). In siena07 the effect has no estimate to compare wi…\\
\opt{settingOppActSqrtD} & one-mode network & Not implemented: RSiena's settings model is a separate model family (setting steps with rates per setting). In siena07 the effect has no estimate to compare wi…\\
\opt{settingOppActLogD} & one-mode network & Not implemented: RSiena's settings model is a separate model family (setting steps with rates per setting). In siena07 the effect has no estimate to compare wi…\\
\opt{settingLogCreationActD} & one-mode network & Not implemented: RSiena's settings model is a separate model family (setting steps with rates per setting). In siena07 the effect has no estimate to compare wi…\\
\bottomrule
\end{longtable}
\normalsize

\subsection*{Placeholder rows of RSiena's effects table (2 rows)}
\footnotesize
\begin{longtable}{@{}p{3.6cm}p{4.0cm}p{6.4cm}@{}}
\toprule
RSiena name & Network / context & Note\\
\midrule
\endhead
\opt{minXAlt} & non-directed network + ?yyyyyy + ?zzzzzz & Placeholder entry of the RSiena effects table; the effect is not offered by siena07 for these variables.\\
\opt{maxXAlt} & non-directed network + ?yyyyyy + ?zzzzzz & Placeholder entry of the RSiena effects table; the effect is not offered by siena07 for these variables.\\
\bottomrule
\end{longtable}
\normalsize

\subsection*{Cross-network effects that also involve a behavior variable (needs behavior inside nwsaom multinet) (1 row)}
\footnotesize
\begin{longtable}{@{}p{3.6cm}p{4.0cm}p{6.4cm}@{}}
\toprule
RSiena name & Network / context & Note\\
\midrule
\endhead
\opt{sameWXClosure} & two-mode network + another network + a behavior variable & Not implemented in nwsaom multinet: a two-mode member together with a co-evolving behavior.\\
\bottomrule
\end{longtable}
\normalsize

\subsection*{Rate effects of continuous behavior (not planned) (1 row)}
\footnotesize
\begin{longtable}{@{}p{3.6cm}p{4.0cm}p{6.4cm}@{}}
\toprule
RSiena name & Network / context & Note\\
\midrule
\endhead
\opt{RateX} & rate (one-mode network) + continuous variable & Not implemented: continuous (stochastic differential equation) behavior variables are not planned.\\
\bottomrule
\end{longtable}
\normalsize

\subsection*{Deliberately out of scope (80 rows)}
\noindent Continuous (SDE) dependent variables and their effects are not part of nwsaom (80 rows of the RSiena effects table).


